% Business Model Canvas. Una página horizontal dentro de modelo_negocio.tex, que
% lo inserta con % Business Model Canvas. Una página horizontal dentro de modelo_negocio.tex, que
% lo inserta con % Business Model Canvas. Una página horizontal dentro de modelo_negocio.tex, que
% lo inserta con % Business Model Canvas. Una página horizontal dentro de modelo_negocio.tex, que
% lo inserta con \input{canvas}. Este archivo no se compila solo.
% Llenar los nueve bloques. Cada afirmación termina con el identificador de la
% hipótesis que la respalda en hipotesis.csv, por ejemplo (H1).
% Marcar las hipótesis validadas con \val{H1} y las invalidadas con \inval{H2}.
\definecolor{navy}{HTML}{1F3467}
\providecommand{\val}[1]{\textcolor{green!50!black}{\textbf{(#1 \checkmark)}}}
\providecommand{\inval}[1]{\textcolor{red!70!black}{\textbf{(#1 $\times$)}}}

% ---------- Contenido de los bloques ----------
\newcommand{\Socios}{\begin{itemize}
  \item Una dependencia del gobierno de la CDMX acepta entregar datos de incidentes con actualización más frecuente que la de sus datos abiertos. (H9)
\end{itemize}}
\newcommand{\Actividades}{\begin{itemize}
  \item Una comunidad de 1,000 usuarios activos genera reportes suficientes para mantener actualizado el mapa de una alcaldía. (H10)
\end{itemize}}
\newcommand{\Recursos}{\begin{itemize}
  \item El riesgo base de cada calle viene de los robos a transeúnte denunciados ante la Fiscalía de la CDMX, y los usuarios lo actualizan con reportes colaborativos de incidentes o zonas inseguras. \val{H8}
\end{itemize}}
\newcommand{\Propuesta}{\begin{itemize}
  \item Los usuarios perciben como más segura la ruta que sugiere la aplicación que la de su planificador habitual. \val{H3}
  \item La ubicación destacada, que muestra sola en el mapa dos productos o servicios del negocio con foto y precio, aumenta al menos 10\% las visitas de clientes a un negocio. \val{H4}
\end{itemize}}
\newcommand{\Relacion}{\begin{itemize}
  \item Los dueños de negocios aceptan una prueba gratuita de 2,000 impresiones con reporte de visitas antes de pagar. \val{H7}
\end{itemize}}
\newcommand{\Canales}{\begin{itemize}
  \item Los dueños de negocios conocen nuevas herramientas de publicidad mediante publicidad masiva en redes sociales, radio o anuncios en vía pública. \val{H5}
  \item Al menos 3\% de las visitas a la landing page se registran para la versión de prueba. \val{H6}
\end{itemize}}
\newcommand{\Segmentos}{\begin{itemize}
  \item Las personas que se desplazan a pie o en transporte público en la CDMX cambian su ruta por seguridad aunque su planificador habitual les sugiera otra más corta. \val{H1}
  \item Los negocios de comercio y servicios ubicados sobre las rutas peatonales de la CDMX buscan más visitas de personas que pasan a pie. \val{H2}
\end{itemize}}
\newcommand{\Costos}{\begin{itemize}
  \item El servidor y el servicio de mapas para 1,000 usuarios activos cuestan menos de 3,000 pesos al mes. \val{H12}
\end{itemize}}
\newcommand{\Ingresos}{\begin{itemize}
  \item Los negocios pagarían 100 pesos por cada 1,000 impresiones como ubicación destacada. \val{H11}
\end{itemize}}

% ---------- Dibujo. No editar debajo de esta línea ----------
\begingroup
\fontsize{9}{11}\selectfont
\setlist{nosep,leftmargin=1em}
\noindent{\Large\bfseries\color{navy} Business Model Canvas}\hfill
{\large GoIA}\hfill 2.0

\vspace{4pt}
\noindent
\resizebox{\linewidth}{!}{%
\begin{tikzpicture}[x=1cm,y=1cm]
  \def\W{25.8}\def\H{16.8}\def\h{5.0}
  \tikzset{caja/.style={draw=navy,line width=0.8pt},
           tit/.style={anchor=north west,font=\bfseries\color{navy}},
           cue/.style={anchor=north west,text width=#1cm,align=left}}
  % Retícula: 5 columnas arriba, 2 abajo
  \def\c{5.16}
  \draw[caja] (0,0) rectangle (\W,\H);
  \draw[caja] (0,\h) -- (\W,\h);
  \foreach \i in {1,2,3,4} \draw[caja] (\i*\c,\h) -- (\i*\c,\H);
  \draw[caja] (\W/2,0) -- (\W/2,\h);
  \pgfmathsetmacro\m{(\h+\H)/2}
  \draw[caja] (\c,\m) -- (2*\c,\m);
  \draw[caja] (3*\c,\m) -- (4*\c,\m);
  % Títulos y contenido
  \node[tit] at (0.1,\H-0.1) {Socios clave};
  \node[cue=4.9] at (0.1,\H-0.7) {\Socios};
  \node[tit] at (\c+0.1,\H-0.1) {Actividades clave};
  \node[cue=4.9] at (\c+0.1,\H-0.7) {\Actividades};
  \node[tit] at (\c+0.1,\m-0.1) {Recursos clave};
  \node[cue=4.9] at (\c+0.1,\m-0.7) {\Recursos};
  \node[tit] at (2*\c+0.1,\H-0.1) {Propuesta de valor};
  \node[cue=4.9] at (2*\c+0.1,\H-0.7) {\Propuesta};
  \node[tit] at (3*\c+0.1,\H-0.1) {Relación con clientes};
  \node[cue=4.9] at (3*\c+0.1,\H-0.7) {\Relacion};
  \node[tit] at (3*\c+0.1,\m-0.1) {Canales};
  \node[cue=4.9] at (3*\c+0.1,\m-0.7) {\Canales};
  \node[tit] at (4*\c+0.1,\H-0.1) {Segmentos de clientes};
  \node[cue=4.9] at (4*\c+0.1,\H-0.7) {\Segmentos};
  \node[tit] at (0.1,\h-0.1) {Estructura de costos};
  \node[cue=12.6] at (0.1,\h-0.7) {\Costos};
  \node[tit] at (\W/2+0.1,\h-0.1) {Fuentes de ingresos};
  \node[cue=12.6] at (\W/2+0.1,\h-0.7) {\Ingresos};
\end{tikzpicture}}

\vspace{2pt}
\noindent\scriptsize Basado en Osterwalder y Pigneur, \emph{Business Model Generation} (2010), CC BY-SA 3.0.
\val{H1} validada,\ \inval{H2} invalidada, (H3) pendiente. Evidencia en \texttt{hipotesis.csv} y \texttt{entrevistas.csv}.
\par\endgroup
. Este archivo no se compila solo.
% Llenar los nueve bloques. Cada afirmación termina con el identificador de la
% hipótesis que la respalda en hipotesis.csv, por ejemplo (H1).
% Marcar las hipótesis validadas con \val{H1} y las invalidadas con \inval{H2}.
\definecolor{navy}{HTML}{1F3467}
\providecommand{\val}[1]{\textcolor{green!50!black}{\textbf{(#1 \checkmark)}}}
\providecommand{\inval}[1]{\textcolor{red!70!black}{\textbf{(#1 $\times$)}}}

% ---------- Contenido de los bloques ----------
\newcommand{\Socios}{\begin{itemize}
  \item Una dependencia del gobierno de la CDMX acepta entregar datos de incidentes con actualización más frecuente que la de sus datos abiertos. (H9)
\end{itemize}}
\newcommand{\Actividades}{\begin{itemize}
  \item Una comunidad de 1,000 usuarios activos genera reportes suficientes para mantener actualizado el mapa de una alcaldía. (H10)
\end{itemize}}
\newcommand{\Recursos}{\begin{itemize}
  \item El riesgo base de cada calle viene de los robos a transeúnte denunciados ante la Fiscalía de la CDMX, y los usuarios lo actualizan con reportes colaborativos de incidentes o zonas inseguras. \val{H8}
\end{itemize}}
\newcommand{\Propuesta}{\begin{itemize}
  \item Los usuarios perciben como más segura la ruta que sugiere la aplicación que la de su planificador habitual. \val{H3}
  \item La ubicación destacada, que muestra sola en el mapa dos productos o servicios del negocio con foto y precio, aumenta al menos 10\% las visitas de clientes a un negocio. \val{H4}
\end{itemize}}
\newcommand{\Relacion}{\begin{itemize}
  \item Los dueños de negocios aceptan una prueba gratuita de 2,000 impresiones con reporte de visitas antes de pagar. \val{H7}
\end{itemize}}
\newcommand{\Canales}{\begin{itemize}
  \item Los dueños de negocios conocen nuevas herramientas de publicidad mediante publicidad masiva en redes sociales, radio o anuncios en vía pública. \val{H5}
  \item Al menos 3\% de las visitas a la landing page se registran para la versión de prueba. \val{H6}
\end{itemize}}
\newcommand{\Segmentos}{\begin{itemize}
  \item Las personas que se desplazan a pie o en transporte público en la CDMX cambian su ruta por seguridad aunque su planificador habitual les sugiera otra más corta. \val{H1}
  \item Los negocios de comercio y servicios ubicados sobre las rutas peatonales de la CDMX buscan más visitas de personas que pasan a pie. \val{H2}
\end{itemize}}
\newcommand{\Costos}{\begin{itemize}
  \item El servidor y el servicio de mapas para 1,000 usuarios activos cuestan menos de 3,000 pesos al mes. \val{H12}
\end{itemize}}
\newcommand{\Ingresos}{\begin{itemize}
  \item Los negocios pagarían 100 pesos por cada 1,000 impresiones como ubicación destacada. \val{H11}
\end{itemize}}

% ---------- Dibujo. No editar debajo de esta línea ----------
\begingroup
\fontsize{9}{11}\selectfont
\setlist{nosep,leftmargin=1em}
\noindent{\Large\bfseries\color{navy} Business Model Canvas}\hfill
{\large GoIA}\hfill 2.0

\vspace{4pt}
\noindent
\resizebox{\linewidth}{!}{%
\begin{tikzpicture}[x=1cm,y=1cm]
  \def\W{25.8}\def\H{16.8}\def\h{5.0}
  \tikzset{caja/.style={draw=navy,line width=0.8pt},
           tit/.style={anchor=north west,font=\bfseries\color{navy}},
           cue/.style={anchor=north west,text width=#1cm,align=left}}
  % Retícula: 5 columnas arriba, 2 abajo
  \def\c{5.16}
  \draw[caja] (0,0) rectangle (\W,\H);
  \draw[caja] (0,\h) -- (\W,\h);
  \foreach \i in {1,2,3,4} \draw[caja] (\i*\c,\h) -- (\i*\c,\H);
  \draw[caja] (\W/2,0) -- (\W/2,\h);
  \pgfmathsetmacro\m{(\h+\H)/2}
  \draw[caja] (\c,\m) -- (2*\c,\m);
  \draw[caja] (3*\c,\m) -- (4*\c,\m);
  % Títulos y contenido
  \node[tit] at (0.1,\H-0.1) {Socios clave};
  \node[cue=4.9] at (0.1,\H-0.7) {\Socios};
  \node[tit] at (\c+0.1,\H-0.1) {Actividades clave};
  \node[cue=4.9] at (\c+0.1,\H-0.7) {\Actividades};
  \node[tit] at (\c+0.1,\m-0.1) {Recursos clave};
  \node[cue=4.9] at (\c+0.1,\m-0.7) {\Recursos};
  \node[tit] at (2*\c+0.1,\H-0.1) {Propuesta de valor};
  \node[cue=4.9] at (2*\c+0.1,\H-0.7) {\Propuesta};
  \node[tit] at (3*\c+0.1,\H-0.1) {Relación con clientes};
  \node[cue=4.9] at (3*\c+0.1,\H-0.7) {\Relacion};
  \node[tit] at (3*\c+0.1,\m-0.1) {Canales};
  \node[cue=4.9] at (3*\c+0.1,\m-0.7) {\Canales};
  \node[tit] at (4*\c+0.1,\H-0.1) {Segmentos de clientes};
  \node[cue=4.9] at (4*\c+0.1,\H-0.7) {\Segmentos};
  \node[tit] at (0.1,\h-0.1) {Estructura de costos};
  \node[cue=12.6] at (0.1,\h-0.7) {\Costos};
  \node[tit] at (\W/2+0.1,\h-0.1) {Fuentes de ingresos};
  \node[cue=12.6] at (\W/2+0.1,\h-0.7) {\Ingresos};
\end{tikzpicture}}

\vspace{2pt}
\noindent\scriptsize Basado en Osterwalder y Pigneur, \emph{Business Model Generation} (2010), CC BY-SA 3.0.
\val{H1} validada,\ \inval{H2} invalidada, (H3) pendiente. Evidencia en \texttt{hipotesis.csv} y \texttt{entrevistas.csv}.
\par\endgroup
. Este archivo no se compila solo.
% Llenar los nueve bloques. Cada afirmación termina con el identificador de la
% hipótesis que la respalda en hipotesis.csv, por ejemplo (H1).
% Marcar las hipótesis validadas con \val{H1} y las invalidadas con \inval{H2}.
\definecolor{navy}{HTML}{1F3467}
\providecommand{\val}[1]{\textcolor{green!50!black}{\textbf{(#1 \checkmark)}}}
\providecommand{\inval}[1]{\textcolor{red!70!black}{\textbf{(#1 $\times$)}}}

% ---------- Contenido de los bloques ----------
\newcommand{\Socios}{\begin{itemize}
  \item Una dependencia del gobierno de la CDMX acepta entregar datos de incidentes con actualización más frecuente que la de sus datos abiertos. (H9)
\end{itemize}}
\newcommand{\Actividades}{\begin{itemize}
  \item Una comunidad de 1,000 usuarios activos genera reportes suficientes para mantener actualizado el mapa de una alcaldía. (H10)
\end{itemize}}
\newcommand{\Recursos}{\begin{itemize}
  \item El riesgo base de cada calle viene de los robos a transeúnte denunciados ante la Fiscalía de la CDMX, y los usuarios lo actualizan con reportes colaborativos de incidentes o zonas inseguras. \val{H8}
\end{itemize}}
\newcommand{\Propuesta}{\begin{itemize}
  \item Los usuarios perciben como más segura la ruta que sugiere la aplicación que la de su planificador habitual. \val{H3}
  \item La ubicación destacada, que muestra sola en el mapa dos productos o servicios del negocio con foto y precio, aumenta al menos 10\% las visitas de clientes a un negocio. \val{H4}
\end{itemize}}
\newcommand{\Relacion}{\begin{itemize}
  \item Los dueños de negocios aceptan una prueba gratuita de 2,000 impresiones con reporte de visitas antes de pagar. \val{H7}
\end{itemize}}
\newcommand{\Canales}{\begin{itemize}
  \item Los dueños de negocios conocen nuevas herramientas de publicidad mediante publicidad masiva en redes sociales, radio o anuncios en vía pública. \val{H5}
  \item Al menos 3\% de las visitas a la landing page se registran para la versión de prueba. \val{H6}
\end{itemize}}
\newcommand{\Segmentos}{\begin{itemize}
  \item Las personas que se desplazan a pie o en transporte público en la CDMX cambian su ruta por seguridad aunque su planificador habitual les sugiera otra más corta. \val{H1}
  \item Los negocios de comercio y servicios ubicados sobre las rutas peatonales de la CDMX buscan más visitas de personas que pasan a pie. \val{H2}
\end{itemize}}
\newcommand{\Costos}{\begin{itemize}
  \item El servidor y el servicio de mapas para 1,000 usuarios activos cuestan menos de 3,000 pesos al mes. \val{H12}
\end{itemize}}
\newcommand{\Ingresos}{\begin{itemize}
  \item Los negocios pagarían 100 pesos por cada 1,000 impresiones como ubicación destacada. \val{H11}
\end{itemize}}

% ---------- Dibujo. No editar debajo de esta línea ----------
\begingroup
\fontsize{9}{11}\selectfont
\setlist{nosep,leftmargin=1em}
\noindent{\Large\bfseries\color{navy} Business Model Canvas}\hfill
{\large GoIA}\hfill 2.0

\vspace{4pt}
\noindent
\resizebox{\linewidth}{!}{%
\begin{tikzpicture}[x=1cm,y=1cm]
  \def\W{25.8}\def\H{16.8}\def\h{5.0}
  \tikzset{caja/.style={draw=navy,line width=0.8pt},
           tit/.style={anchor=north west,font=\bfseries\color{navy}},
           cue/.style={anchor=north west,text width=#1cm,align=left}}
  % Retícula: 5 columnas arriba, 2 abajo
  \def\c{5.16}
  \draw[caja] (0,0) rectangle (\W,\H);
  \draw[caja] (0,\h) -- (\W,\h);
  \foreach \i in {1,2,3,4} \draw[caja] (\i*\c,\h) -- (\i*\c,\H);
  \draw[caja] (\W/2,0) -- (\W/2,\h);
  \pgfmathsetmacro\m{(\h+\H)/2}
  \draw[caja] (\c,\m) -- (2*\c,\m);
  \draw[caja] (3*\c,\m) -- (4*\c,\m);
  % Títulos y contenido
  \node[tit] at (0.1,\H-0.1) {Socios clave};
  \node[cue=4.9] at (0.1,\H-0.7) {\Socios};
  \node[tit] at (\c+0.1,\H-0.1) {Actividades clave};
  \node[cue=4.9] at (\c+0.1,\H-0.7) {\Actividades};
  \node[tit] at (\c+0.1,\m-0.1) {Recursos clave};
  \node[cue=4.9] at (\c+0.1,\m-0.7) {\Recursos};
  \node[tit] at (2*\c+0.1,\H-0.1) {Propuesta de valor};
  \node[cue=4.9] at (2*\c+0.1,\H-0.7) {\Propuesta};
  \node[tit] at (3*\c+0.1,\H-0.1) {Relación con clientes};
  \node[cue=4.9] at (3*\c+0.1,\H-0.7) {\Relacion};
  \node[tit] at (3*\c+0.1,\m-0.1) {Canales};
  \node[cue=4.9] at (3*\c+0.1,\m-0.7) {\Canales};
  \node[tit] at (4*\c+0.1,\H-0.1) {Segmentos de clientes};
  \node[cue=4.9] at (4*\c+0.1,\H-0.7) {\Segmentos};
  \node[tit] at (0.1,\h-0.1) {Estructura de costos};
  \node[cue=12.6] at (0.1,\h-0.7) {\Costos};
  \node[tit] at (\W/2+0.1,\h-0.1) {Fuentes de ingresos};
  \node[cue=12.6] at (\W/2+0.1,\h-0.7) {\Ingresos};
\end{tikzpicture}}

\vspace{2pt}
\noindent\scriptsize Basado en Osterwalder y Pigneur, \emph{Business Model Generation} (2010), CC BY-SA 3.0.
\val{H1} validada,\ \inval{H2} invalidada, (H3) pendiente. Evidencia en \texttt{hipotesis.csv} y \texttt{entrevistas.csv}.
\par\endgroup
. Este archivo no se compila solo.
% Llenar los nueve bloques. Cada afirmación termina con el identificador de la
% hipótesis que la respalda en hipotesis.csv, por ejemplo (H1).
% Marcar las hipótesis validadas con \val{H1} y las invalidadas con \inval{H2}.
\definecolor{navy}{HTML}{1F3467}
\providecommand{\val}[1]{\textcolor{green!50!black}{\textbf{(#1 \checkmark)}}}
\providecommand{\inval}[1]{\textcolor{red!70!black}{\textbf{(#1 $\times$)}}}

% ---------- Contenido de los bloques ----------
\newcommand{\Socios}{\begin{itemize}
  \item Una dependencia del gobierno de la CDMX acepta entregar datos de incidentes con actualización más frecuente que la de sus datos abiertos. (H9)
\end{itemize}}
\newcommand{\Actividades}{\begin{itemize}
  \item Una comunidad de 1,000 usuarios activos genera reportes suficientes para mantener actualizado el mapa de una alcaldía. (H10)
\end{itemize}}
\newcommand{\Recursos}{\begin{itemize}
  \item El riesgo base de cada calle viene de los robos a transeúnte denunciados ante la Fiscalía de la CDMX, y los usuarios lo actualizan con reportes colaborativos de incidentes o zonas inseguras. \val{H8}
\end{itemize}}
\newcommand{\Propuesta}{\begin{itemize}
  \item Los usuarios perciben como más segura la ruta que sugiere la aplicación que la de su planificador habitual. \val{H3}
  \item La ubicación destacada, que muestra sola en el mapa dos productos o servicios del negocio con foto y precio, aumenta al menos 10\% las visitas de clientes a un negocio. \val{H4}
\end{itemize}}
\newcommand{\Relacion}{\begin{itemize}
  \item Los dueños de negocios aceptan una prueba gratuita de 2,000 impresiones con reporte de visitas antes de pagar. \val{H7}
\end{itemize}}
\newcommand{\Canales}{\begin{itemize}
  \item Los dueños de negocios conocen nuevas herramientas de publicidad mediante publicidad masiva en redes sociales, radio o anuncios en vía pública. \val{H5}
  \item Al menos 3\% de las visitas a la landing page se registran para la versión de prueba. \val{H6}
\end{itemize}}
\newcommand{\Segmentos}{\begin{itemize}
  \item Las personas que se desplazan a pie o en transporte público en la CDMX cambian su ruta por seguridad aunque su planificador habitual les sugiera otra más corta. \val{H1}
  \item Los negocios de comercio y servicios ubicados sobre las rutas peatonales de la CDMX buscan más visitas de personas que pasan a pie. \val{H2}
\end{itemize}}
\newcommand{\Costos}{\begin{itemize}
  \item El servidor y el servicio de mapas para 1,000 usuarios activos cuestan menos de 3,000 pesos al mes. \val{H12}
\end{itemize}}
\newcommand{\Ingresos}{\begin{itemize}
  \item Los negocios pagarían 100 pesos por cada 1,000 impresiones como ubicación destacada. \val{H11}
\end{itemize}}

% ---------- Dibujo. No editar debajo de esta línea ----------
\begingroup
\fontsize{9}{11}\selectfont
\setlist{nosep,leftmargin=1em}
\noindent{\Large\bfseries\color{navy} Business Model Canvas}\hfill
{\large GoIA}\hfill 2.0

\vspace{4pt}
\noindent
\resizebox{\linewidth}{!}{%
\begin{tikzpicture}[x=1cm,y=1cm]
  \def\W{25.8}\def\H{16.8}\def\h{5.0}
  \tikzset{caja/.style={draw=navy,line width=0.8pt},
           tit/.style={anchor=north west,font=\bfseries\color{navy}},
           cue/.style={anchor=north west,text width=#1cm,align=left}}
  % Retícula: 5 columnas arriba, 2 abajo
  \def\c{5.16}
  \draw[caja] (0,0) rectangle (\W,\H);
  \draw[caja] (0,\h) -- (\W,\h);
  \foreach \i in {1,2,3,4} \draw[caja] (\i*\c,\h) -- (\i*\c,\H);
  \draw[caja] (\W/2,0) -- (\W/2,\h);
  \pgfmathsetmacro\m{(\h+\H)/2}
  \draw[caja] (\c,\m) -- (2*\c,\m);
  \draw[caja] (3*\c,\m) -- (4*\c,\m);
  % Títulos y contenido
  \node[tit] at (0.1,\H-0.1) {Socios clave};
  \node[cue=4.9] at (0.1,\H-0.7) {\Socios};
  \node[tit] at (\c+0.1,\H-0.1) {Actividades clave};
  \node[cue=4.9] at (\c+0.1,\H-0.7) {\Actividades};
  \node[tit] at (\c+0.1,\m-0.1) {Recursos clave};
  \node[cue=4.9] at (\c+0.1,\m-0.7) {\Recursos};
  \node[tit] at (2*\c+0.1,\H-0.1) {Propuesta de valor};
  \node[cue=4.9] at (2*\c+0.1,\H-0.7) {\Propuesta};
  \node[tit] at (3*\c+0.1,\H-0.1) {Relación con clientes};
  \node[cue=4.9] at (3*\c+0.1,\H-0.7) {\Relacion};
  \node[tit] at (3*\c+0.1,\m-0.1) {Canales};
  \node[cue=4.9] at (3*\c+0.1,\m-0.7) {\Canales};
  \node[tit] at (4*\c+0.1,\H-0.1) {Segmentos de clientes};
  \node[cue=4.9] at (4*\c+0.1,\H-0.7) {\Segmentos};
  \node[tit] at (0.1,\h-0.1) {Estructura de costos};
  \node[cue=12.6] at (0.1,\h-0.7) {\Costos};
  \node[tit] at (\W/2+0.1,\h-0.1) {Fuentes de ingresos};
  \node[cue=12.6] at (\W/2+0.1,\h-0.7) {\Ingresos};
\end{tikzpicture}}

\vspace{2pt}
\noindent\scriptsize Basado en Osterwalder y Pigneur, \emph{Business Model Generation} (2010), CC BY-SA 3.0.
\val{H1} validada,\ \inval{H2} invalidada, (H3) pendiente. Evidencia en \texttt{hipotesis.csv} y \texttt{entrevistas.csv}.
\par\endgroup
